Táto práca sa zaoberá modelovaním a riadením doskového tepelného výmenníka pomocou metód prediktívneho riadenia. Cieľom práce bolo identifikovať dynamiku systému z experimentálnych meraní, navrhnúť regulátory typu Model Predictive Control (MPC) s využitím dvoch modelovacích prístupov a porovnať ich správanie v simuláciách a experimentoch.
Experimentálny systém pozostával zo zásobníka, doskového tepelného výmenníka a čerpadiel zabezpečujúcich cirkuláciu kvapaliny v systéme. Ako riadiaci vstup bol použitý výkon čerpadla, ktoré privádzalo horúcu kvapalinu do výmenníka. Regulovanou veličinou bola výstupná teplota z tepelného výmenníka.
Identifikácia systému bola vykonaná pomocou skokových zmien prietoku horúcej kvapaliny. Na základe nameraných dát bol identifikovaný lineárny model prvého rádu. Okrem toho bol vytvorený dátovo orientovaný Koopmanov model, ktorý umožňuje zlepšenú aproximáciu nelineárneho správania systému pomocou lineárnej dynamiky vo viac rozmernom priestore.
Pre oba modely boli navrhnuté MPC regulátory s kvadratickou optimalizačnou úlohou a s obmedzeniami na vstup a výstup systému. Súčasťou riadiacej štruktúry bol Kalmanov filter, ktorý slúžil na odhad stavov systému. Modely boli porovnávané z hľadiska presnosti riadenia, odchýlky od ustáleného stavu a veľkosti regulačného zásahu. Simulácie boli vykonané s využitím nelineárneho matematického modelu zariadenia.
V experimentálnej časti práce bolo MPC riadenie aplikované na tepelný výmenník, pričom systém bol riadený z rôznych počiatočných teplôt do ustáleného stavu. Koopmanov MPC vo väčšine prípadov zlepšil prechodovú odozvu a znížil regulačnú chybu, avšak za cenu agresívnejších regulačných zásahov v porovnaní s MPC založeným na lineárnom modeli prvého rádu.

\vspace{0.5cm}
\noindent\textbf{Kľúčové slová:} doskový výmenník tepla, prediktívne riadenie, Koopmanov model, identifikácia systému